\documentclass[10pt,a4paper]{amsart}
\usepackage[margin=19mm]{geometry}
\usepackage{amsmath,amssymb,mathtools}
\usepackage[scr=rsfs,cal=euler]{mathalfa}
\usepackage{xfrac}
\usepackage[unicode,hidelinks,bookmarks=false]{hyperref}
\newcommand{\ie}{i.e.\ }
\newcommand{\cf}{cf.\ }
\newcommand{\resp}{respectively\ }
\newcommand{\Subsection}{\subsection}
\providecommand{\nobreakdash}{\nobreak\hbox{-}}
\setlength{\emergencystretch}{2em}
\multlinegap0pt
\usepackage{bm}
\usepackage{bbm}
\usepackage{dsfont}
\usepackage{xspace}
\usepackage{cancel}
\usepackage{mathtools}
\usepackage{amsthm}
\makeatletter
\@ifundefined{theorem}{\newtheorem{theorem}{Theorem}}{}
\@ifundefined{proposition}{\newtheorem{proposition}[theorem]{Proposition}}{}
\makeatother
\makeatletter
\expandafter\ifx\csname proof\endcsname\relax
\newenvironment{proof}[1][Proof]{\noindent\textbf{#1.} }{\ \rule{0.5em}{0.5em}}
\fi
\makeatother
\title[Heisenberg uncertainty inequalities for locally compact abel]{Heisenberg uncertainty inequalities for locally compact abelian groups}
\author{Hartmut F\"uhr}
\date{Source: arXiv:2606.30079v1 (CC BY 4.0)}
\begin{document}
\maketitle

\noindent\emph{Source and licence.} Heisenberg uncertainty inequalities for locally compact abelian groups. Version arXiv:2606.30079v1 (CC BY 4.0), distributed under the Creative Commons Attribution 4.0 International licence, as recorded in the arXiv licence metadata for this exact version and shown on the abstract page https://arxiv.org/abs/2606.30079v1. Full rights notice: licenses/arxiv-2606.30079-CC-BY-4.0.txt.\par\smallskip

\noindent\textit{Editorial scope.} Counted unit: the Proposition whose source label is \texttt{prop:main} in the article (source0.tex lines 353--358, source label \texttt{prop:main}), delivered together with the article's own complete original proof of it (source0.tex lines 359--397, one \texttt{proof} environment of 39 lines). No mathematics is written, repaired or re-derived: statement and proof are the article's own text line for line. Required context retained uncounted: none. Every definition, display and label of those lines is retained unchanged. Omitted from this package and not redistributed: 1--352, 398--609. Nothing inside a retained excerpt is omitted or altered: every retained body is delivered exactly as the article's lines are written, so no omission receipt of any kind accompanies this package. Citations: the retained excerpts carry no citation command at all, so no bibliography entry of the article is transcribed and no reference is redistributed. Reference adaptation: none was needed and none was made; every reference inside the delivered unit resolves inside it, so no unresolved reference or citation is produced. Printed slips kept literal: no printed word, symbol, hypothesis or number of the counted statement or of its proof was altered. Layout and class: the article's own class is \texttt{amsart} with packages including amsmath, amsfonts, amssymb, comment, graphicx, xcolor, enumitem, geometry; the wrapper is the lane's pinned \texttt{amsart} editorial wrapper at 10pt on A4, which restates the theorem-like environments the retained text needs and adds the article's own macro definitions that the retained text needs (none), with extra wrapper packages (bm, bbm, dsfont, xspace, cancel, mathtools). The delivered numbering is the unit's own. Assets: no retained span contains a figure environment, a graphics-inclusion command, a PDF-page inclusion, a table or a TikZ picture; no image file is copied into this package, the package declares no asset dependency and no image file is redistributed with it. Licence: version arXiv:2606.30079v1 of this article is distributed under the Creative Commons Attribution 4.0 International licence, as recorded in the arXiv licence metadata for this exact version and shown on the abstract page https://arxiv.org/abs/2606.30079v1. A full rights notice is delivered at licenses/arxiv-2606.30079-CC-BY-4.0.txt. Characters: every retained excerpt is pure ASCII, so the delivered engine, XeLaTeX with its default Latin Modern text font, reports no missing character.
\begin{proposition} \label{prop:main}
Let $\delta \in (0,1/2)$ be given, as well as $f \in L^2(G)$ with $\| f \|_2 = 1$ and $\| Tf \|_2 < \infty$. Then there exists $r_0 >0$ satisfying 
\[
\forall r > r_0: \| P_{r} f \|_2 \ge 1 - \delta~, \mbox{ and } r_0 \le \frac{\| T f \|_2}{\left(2 \delta - \delta^2 \right)^{1/2}}~.
\]
\end{proposition}

\begin{proof} The core idea of the proof is to derive a concentration estimate by applying Markov's inequality to a suitably defined random variable. The measure decomposition over the $d$-spheres allows to write
\begin{eqnarray*}
 1 & = & \int_G |f(x)|^2 dx \\
  & = & \int_{\mathbb{R}_0^+} \int_{S_r(e)} |f(x)|^2 d\mu_r(x) d\nu(r)~,
\end{eqnarray*}
hence letting  
\[
d\nu_f(r) = \int_{S_r(e)} |f(x)|^2 d\mu_r(x) d\nu(r)
\] defines a probability measure $\nu_f(r)$ on $\mathbb{R}_0^+$. From now on let $R$ denote a random variable with probability given by $\nu_f$. We then obtain
\begin{eqnarray}
 \nonumber    \mathbb{E}(R^2) & = & \int_{\mathbb{R}_0^+} r^2 \int_{S_r(e)} |f(x)|^2 d\mu_r(x) d\nu(r) \\
 \nonumber   & = & \int_{\mathbb{R}_0^+} \int_{S_r(e)} d(x,e)^2 |f(x)|^2 d\mu_r(x) d\nu(r) \\
 \label{eqn:expectation}   & = & \int_G d(x,e)^2 |f(x)|^2 dx = \| T f \|_2^2 ~,
\end{eqnarray}
where the last equality used the measure decomposition property. 
Now let 
\[
r_0 = \sup \left\{ r \ge 0: \int_{d(x,e) \ge r} |f(x)|^2 dx > 2 \delta - \delta^2 \right\}~.
\]
$r_0 \in (0,\infty)$ is well-defined since 
\[
\int_{d(x,e) \ge r} |f(x)|^2 dx = \mathbb{P} (R \ge r) \to 0 ~,
\]
 as $r \to \infty$. 
 Now, on the one hand, the definition of $r_0$ as supremum implies for $r>r_0$
 \begin{equation} \label{eqn:deliverable_1}
  \| P_{r} f \|_2^2  = 1 - \int_{d(x,e) \ge r} |f(x)|^2 dx \ge 1-2 \delta + \delta^2 = (1-\delta)^2~. 
 \end{equation}
 
 For any $r < r_0$ one has by definition 
 \[ 2 \delta - \delta^2 \le \int_{d(x,e) \ge r}|f(x)|^2 dx = P(R \ge r) \le \frac{\mathbb{E}(R^2)}{r^2} = \frac{\| T f \|_2^2}{r^2}~, \]
 where we used the Markov inequality for the second moment, as well as equation (\ref{eqn:expectation}).
This yields equivalently 
\[ r^2 \le \frac{\| Tf \|_2^2}{2 \delta - \delta^2}
\] for all $r < r_0$, and then $r \to r_0$ yields
\[
r_0 \le \frac{\| T f\|_2}{\left(2 \delta - \delta^2\right)^{1/2}}~.
\]
\end{proof}

\end{document}
